\section{Adjoint regularity and monodromy}
\label{sec:monodromy}

The semisimple Weil parameter does not determine individual local
$L$-factors, since those factors also involve monodromy.  We first show
that the ratio occurring in a local functional equation is nevertheless
determined up to a unit.  We then explain how purity controls local
$L$-factors and determinants.  Finally, we prove that regularity of the
adjoint $L$-factor at $1$ determines the monodromy orbit.  These statements
will allow us to compare two monodromy operators once their semisimple
Weil parameters have been identified, without assuming compatibility of
the monodromy operators themselves.

\subsection{A ratio independent of monodromy}

\begin{lemma}\label{lem:monodromy-ratio}
Let $r$ be a Frobenius-semisimple Weil representation on a
finite-dimensional complex vector space $V$, with finite inertial
image.  If $N$ and $N'$ are nilpotent operators such that $(r,N)$ and
$(r,N')$ are Weil--Deligne representations, then
\[
 \frac{L(1-s,(r,N)^\vee)}{L(s,(r,N))}
 \doteq
 \frac{L(1-s,(r,N')^\vee)}{L(s,(r,N'))}.
\]
More precisely, put $W=V^{r(I_k)}$, $K=\ker(N|_W)$, and
$F=r(\mathrm{Fr})$.  If $R_N(s)$ denotes the ratio on the left, then
\begin{equation}\label{eq:monodromy-ratio-unit}
 \frac{R_0(s)}{R_N(s)}
   =\det\bigl(-Q^{-s}F\mid W/K\bigr).
\end{equation}
\end{lemma}

\begin{proof}
The operator $N$ commutes with inertia and satisfies
$FN=Q^{-1}NF$.  Thus $N$ identifies $W/K$ with $NW$, multiplying
Frobenius eigenvalues by $Q^{-1}$.  Averaging over the finite inertia
image identifies $(V^\vee)^{r^\vee(I_k)}$ with $W^\vee$.
The dual monodromy operator is $-N^{\mathsf t}$.  On $W^\vee$ its
kernel is the annihilator of $NW$, so restriction of functionals gives
the Frobenius-equivariant isomorphism
\[
 W^\vee/\ker(-N^{\mathsf t})\simeq (NW)^\vee.
\]
Consequently, if $c$ is a Frobenius eigenvalue on $W/K$, the
corresponding eigenvalue on this dual quotient is $Qc^{-1}$.
Writing $x=Q^{-s}$ and comparing the determinants for $N$ and $0$
gives
\[
 \frac{R_0(s)}{R_N(s)}
   =\prod_c\frac{1-xc}{1-Q^{-1}x^{-1}(Qc^{-1})}
   =\prod_c(-xc),
\]
where the eigenvalues are counted with multiplicity.  This proves
\eqref{eq:monodromy-ratio-unit}, as an identity of rational functions
in $x$.  Its right-hand side is a nonzero constant times a monomial,
and comparison through $N=0$ proves the assertion for $N$ and $N'$.
\end{proof}

\subsection{Consequences of purity}

Recall that the \emph{monodromy filtration centered at zero} of a
nilpotent operator $N$ on $V$ is the unique finite increasing
filtration $M_\bullet V$, indexed by the integers, satisfying
\[
 N(M_jV)\subset M_{j-2}V,
 \qquad
 N^j:\operatorname{Gr}^{M}_jV
      \xrightarrow{\ \sim\ }
      \operatorname{Gr}^{M}_{-j}V
 \quad (j\geq 0).
\]
Here $M_jV=0$ for sufficiently negative $j$ and $M_jV=V$ for
sufficiently positive $j$.  On a Jordan chain
$v,Nv,\ldots,N^dv$, the successive vectors have degrees
$d,d-2,\ldots,-d$.  This description constructs the filtration and
shows that it commutes with direct sums.  Its uniqueness also shows
that the Weil action preserves it, since conjugation by $r(w)$
multiplies $N$ by the nonzero scalar $|w|$.

We say that $(r,N)$ is \emph{pure of weight zero}, with respect to the
fixed complex realization of the coefficients, if every Frobenius
eigenvalue on $\operatorname{Gr}^{M}_jV$ has absolute value
$Q^{j/2}$.  When $N=0$, this says that all eigenvalues of
$r(\mathrm{Fr})$ have absolute value $1$.  Nonzero monodromy allows
eigenvalues of different absolute values, arranged in symmetric
Jordan chains.

\begin{lemma}\label{lem:pure-determinant}
Let $(r,N)$ be a finite-dimensional Frobenius-semisimple
Weil--Deligne representation, pure of weight zero.  Then
$L(s,(r,N))$ is holomorphic and nonzero for $\operatorname{Re}(s)>0$,
and
\begin{equation}\label{eq:pure-determinant}
 \left|\det r(\mathrm{Fr})\right|=1.
\end{equation}
Every Weil--Deligne direct summand of $(r,N)$ is also pure of weight
zero and satisfies these conclusions.
\end{lemma}

\begin{proof}
The bottom vector of a Jordan chain of length $d+1$ has monodromy
degree $-d$.  Hence $\ker N\subset M_0V$.  Purity implies that all
Frobenius eigenvalues on $M_0V$, and therefore on $(\ker N)^{I_k}$,
have absolute value at most $1$.  If $\operatorname{Re}(s)>0$, each
factor $1-Q^{-s}c$ in the defining determinant of the local
$L$-factor is nonzero.  The asserted holomorphy and nonvanishing
follow.

The defining isomorphisms of the monodromy filtration give
$\dim\operatorname{Gr}^{M}_jV
 =\dim\operatorname{Gr}^{M}_{-j}V$.  Taking the determinant on the
graded spaces therefore gives
\[
 \left|\det r(\mathrm{Fr})\right|
   =Q^{\frac12\sum_j j\dim\operatorname{Gr}^{M}_jV}=1.
\]
Finally, the monodromy filtration of a direct sum is the direct sum
of its monodromy filtrations.  Each graded space of a
Weil--Deligne summand is consequently a Frobenius-stable direct
summand of the corresponding graded space of $V$, and inherits its
purity.
\end{proof}

For example, take trivial inertia on $V=\mathbb C e_1\oplus
\mathbb C e_2$ and set
\[
 r(\mathrm{Fr})=
 \begin{pmatrix}Q^{-1/2}&0\\0&Q^{1/2}\end{pmatrix},
 \qquad
 Ne_2=e_1,\quad Ne_1=0.
\]
The degrees of $e_1,e_2$ are $-1,1$, so this representation is pure
of weight zero although the image of $r$ is unbounded.  It comes
from the $L$-parameter trivial on $W_k$ and standard on
$\mathrm{SL}_2(\mathbb C)$.  Replacing $N$ by $0$ destroys purity
and introduces a pole at $s=1/2$ in the local $L$-factor.  This is
why boundedness of the Weil part and purity of a completed
Weil--Deligne representation must be distinguished.

\subsection{Adjoint regularity determines the orbit}

The next proposition is the implication from adjoint regularity to an
open monodromy orbit established in
\cite[Proposition~3.5]{CDFZ2025} and
\cite[Proposition~6.10]{DHKM2026}.  We give the elementary argument in
the semisimple case needed here.  The Killing form identifies the
obstruction to an orbit being open with the space responsible for a
pole of the adjoint $L$-factor at~$1$.

\begin{proposition}\label{prop:open-orbit}
Let $\widehat G$ be a connected semisimple complex algebraic group,
and let $r:W_k\to\widehat G$ be a Frobenius-semisimple Weil
parameter with finite inertial image.  Write
\[
 \mathfrak e=\operatorname{Lie}(\widehat G)^{r(I_k)},
 \qquad
 \mathfrak e_a=
 \ker\bigl(\operatorname{Ad}(r(\mathrm{Fr}))-a
                 \mid\mathfrak e\bigr)
 \quad(a\in\mathbb C^\times),
\]
and let $C=Z_{\widehat G}(r(W_k))$.  If $(r,N)$ is a
Weil--Deligne parameter and $L(s,\operatorname{Ad}\circ(r,N))$
is regular at $s=1$, then the orbit $C\cdot N$ is open in
$\mathfrak e_{Q^{-1}}$.
In particular, if $N_1,N_2$ are two monodromy operators over $r$
whose adjoint $L$-factors are regular at $1$, then some $c\in C$
satisfies $\operatorname{Ad}(c)N_1=N_2$.
\end{proposition}

\begin{proof}
The Weil--Deligne relation places $N$ in
$\mathfrak e_{Q^{-1}}$.  On the adjoint representation its
monodromy is $\operatorname{ad}(N)$, so regularity at $1$ is
equivalent to
\begin{equation}\label{eq:adjoint-pole-kernel}
 \ker(\operatorname{ad}N)\cap\mathfrak e_Q=0.
\end{equation}
Indeed, a pole at $1$ means that Frobenius has eigenvalue $Q$ on
$(\ker\operatorname{ad}N)^{r(I_k)}$.

Let $B$ be the Killing form of $\operatorname{Lie}(\widehat G)$.
Its restriction to $\mathfrak e$ is nondegenerate.  To see this,
average any vector over the finite inertia image: pairing it with
an inertia-invariant vector gives the same value before and after
averaging.  A vector orthogonal to $\mathfrak e$ inside
$\mathfrak e$ is therefore orthogonal to the entire Lie algebra,
and is zero.  Frobenius preserves $B$ and acts semisimply, so $B$
pairs $\mathfrak e_a$ perfectly with $\mathfrak e_{a^{-1}}$.

We have $\operatorname{Lie}(C)=\mathfrak e_1$.  The differential
of the orbit map at the identity is
\begin{equation}\label{eq:monodromy-orbit-tangent}
 \mathfrak e_1\longrightarrow\mathfrak e_{Q^{-1}},
 \qquad Y\longmapsto[Y,N].
\end{equation}
By invariance of $B$, its transpose under the preceding perfect
pairings is
\[
 \mathfrak e_Q\longrightarrow\mathfrak e_1,
 \qquad X\longmapsto[N,X].
\]
The kernel of this transpose is zero by
\eqref{eq:adjoint-pole-kernel}, so
\eqref{eq:monodromy-orbit-tangent} is surjective.  Algebraic group
orbits are smooth and locally closed in characteristic zero.
Thus $C\cdot N$ has the dimension of $\mathfrak e_{Q^{-1}}$ and
is open in that vector space.  Since the vector space is
irreducible, it has at most one open $C$-orbit.  This proves the
last assertion as well.
\end{proof}
